\scanpage{142}
\mspage{V.14}
% archive sheet number: 141
% manuscript lines 1-8
\uline{Remarque} 3.2. a) Lorsque $X_s$ est \inserted{localement} une intersection complète, \deleted{alors} et purement de dimension $n$, alors \deleted{\unsure{l'hypothèse}} \deleted{les hypothèses} il en est de même de $X_{\unsure{\eta|y}}$, et les hypothèses (3.1.1) et (3.1.2) sont vérifiées, avec $m=n$. Nous verrons cependant plus bas (au \deleted{\unsure{\S}}\S\ 5) que l'on peut alors (sous réserve de conjectures rappelées \deleted{dans (\unsure{2.2})} tout à l'heure 3.3)) améliorer le résultat \inserted{(3.1.3)} en écrivant $\underline\Phi^i_{\unsure{\nu|v}}=0$ pour $i\leq n-d$.

% manuscript lines 9-13
b) Par contre, \deleted{\unsure{bien}} même si $X_{\bar\eta}$ et $X_s$ sont \deleted{\unsure{irréductibles}} \inserted{purement} de dimension $n$, \deleted{et si $n=1$}, et si les hypothèses (3.1.1) et (3.1.2) sont satisfaites \deleted{pour} avec $m=n$, on ne peut en général améliorer la conclusion (3.1.3) d'une unité comme dans a).

% manuscript lines 14-23
\deleted{Pour} Pour le voir, \unsure{partons|prenons} p. ex. de \deleted{\unsure{un} $X'$ lisse sur $S$ à fibres géom. \deleted{\unsure{int.}} connexes (p. ex.} \deleted{$X'=$} $X'=\mathbb P^1_S\amalg\mathbb P^1_S$ et \deleted{\unsure{par}} un $X$ déduit de $X'$ en identifiant deux pts \inserted{distincts} de $X'_s$ rationnels sur $k$, \inserted{appartenant à deux composantes distinctes de $X'_s$}. Alors \deleted{les hypothèses} $X$ satisfait les hypothèses avec $m=n=1$, $d=0$, cependant \deleted{la conclusion} on n'a pas $\Phi^1=0$, \deleted{car \unsure{contrairement}} \inserted{i.e.} l'hom. canonique $H^0(X_0)\mathbin{-}H^{\unsure{0|1}}(X_{\bar\eta})$ n'est pas bijectif, puisque $X_0$ est connexe mais $X_{\bar\eta}$ ne l'est pas.
